\documentclass[10pt,a4paper]{amsart}
\usepackage[margin=19mm]{geometry}
\usepackage{amsmath,amssymb,mathtools}
\usepackage[scr=rsfs,cal=euler]{mathalfa}
\usepackage{xfrac}
\usepackage[unicode,hidelinks,bookmarks=false]{hyperref}
\newcommand{\ie}{i.e.\ }
\newcommand{\cf}{cf.\ }
\newcommand{\resp}{respectively\ }
\newcommand{\Subsection}{\subsection}
\providecommand{\nobreakdash}{\nobreak\hbox{-}}
\setlength{\emergencystretch}{2em}
\multlinegap0pt
\usepackage{xcolor}
\usepackage{amssymb}
\usepackage[shortlabels]{enumitem}
\usepackage{mathtools}
\usepackage{siunitx}
\usepackage{bm}
\usepackage{tikz}
\newtheorem{lemma}{Lemma}
\newtheorem{theorem}{Theorem}
\newcommand{\Z}{\mathbb{Z}}
\newcommand{\sabnam}[1]{{\color{red} #1}}
\title{On the Schur multiplier of $p$-groups with abelianization $s$-elementary abelian}
\author{Sumana Hatui, Tony Nixon Mavely, Sahanawaj Sabnam}
\date{Source: arXiv:2605.00810v1 (CC BY 4.0)}
\begin{document}
\maketitle

\noindent\emph{Source and licence.} On the Schur multiplier of $p$-groups with abelianization $s$-elementary abelian. Version arXiv:2605.00810v1 (CC BY 4.0), distributed by arXiv under the Creative Commons Attribution 4.0 International licence, http://creativecommons.org/licenses/by/4.0/, as recorded in the arXiv licence metadata for this exact version and shown on the abstract page https://arxiv.org/abs/2605.00810v1. Full licence notice: \texttt{licenses/arxiv-2605.00810-CC-BY-4.0.txt}.\par\smallskip

\noindent\textit{Editorial scope.} Counted unit: the article's theorem M(G\_(j,k (source0.tex lines 1006--1021). The article's complete original proof of it is delivered with it (source0.tex lines 1023--1109, one \texttt{proof} environment of 87 lines). No mathematics is written, repaired or re-derived: the statement and the proof are the article's own lines, in the article's own notation, and no retained line is substituted or re-flowed.  Retained uncounted as the source-local context the proof needs: lines 794--821.  Nothing was omitted from any retained span, so the package carries no omission record.  No retained line contains a citation, so the wrapper carries no bibliography and no bibliography entry.  No retained line contains a figure environment, an image-inclusion command or drawing code, so no image is delivered and the package declares no asset dependency.  Editorial formatting only, in addition: an AMS \texttt{amsart} preamble that restates the article's own declarations and macros used by the retained lines, namely the environments \texttt{lemma}, \texttt{theorem} on separate counters, together with the standard packages those lines need.  The retained environments print in this document as Lemma 1, Theorem 1 in order of appearance, whereas the article prints the same items with the numbering of its own full document.  The complete native source of the article as distributed in the arXiv e-print for this exact version is bundled unchanged as \texttt{sources/199594/source0.tex}.
\begin{lemma}  \label{SchurG_{j,k}} 
Let $d,j,k \in \mathbb{N}$ such that $d \geq 2$, $0 \leq k \leq d-2$, ${1 \leq j \leq k+1}$.
   Consider a group $G_{j,k}^d$, minimally generated by $\alpha_1,\alpha_2,\ldots,\alpha_d$, having the following presentation,
%\[G_{j,k}^d=\Biggl\langle 
%\begin{array}{l|cl}
%	&  [\alpha_i,\alpha_{i+1}]=\beta_i, 1 \leq i \leq k+1,   %\\
%	 \alpha_i, 1 \le i \le d&  \alpha_i^{p^s}=\beta_i, \beta_i^{p^{t_i}}=1 \text{ for } 1 \leq i \leq j, \\
%     & \alpha_r^{p^s}= \beta_l^{p^s}=1 \text{ for } j+1 \leq r \leq d  \text{ and } j+1 \leq l \leq k+1
%\end{array}
%\Biggr\rangle\]
\[G_{j,k}^d=\Biggl\langle 
\begin{array}{l|cl}
\alpha_i, 1 \le i \le d,	&  [\alpha_i,\alpha_{i+1}]=\beta_i, 1 \leq i \leq k+1,   \\
	\beta_u,1 \le u \le k+1 &  \alpha_i^{p^s}=\beta_i, \alpha_l^{p^s}=1 \text{ for } 1 \leq i \leq j <l \leq d, \\
     &  \beta_i^{p^{t_i}}= \beta_l^{p^s}=1 \text{ for }  1 \leq i \leq j < l \leq k+1
\end{array}
\Biggr\rangle.\]
    % $G_{j,k}^d=\langle \alpha_i,1 \le i \le d \mid [\alpha_1,\alpha_2]=\beta_1,[\alpha_2,\alpha_3]=\beta_2,\ldots , [\alpha_{k+1},\alpha_{k+2}]=\beta_{k+1},\alpha_1^{p^s}=\beta_1, \alpha_2^{p^s}=\beta_2,\ldots,\alpha_j^{p^s}=\beta_j ,\beta_1^{p^{t_1}}=\beta_2^{p^{t_2}}=\cdots=\beta_j^{p^{t_j}}=1, \beta_l^{p^s}=1,\alpha_r^{p^s}=1,\text{ for } 0 \le k \le d-2, 1 \le j \le k+1, j+1 \le l \le { k+1}, j+1 \le r \le d \rangle.$\\
\begin{enumerate} 
    \item [(i)]  If $j < k+1$, then 
%$$M(G_{j,k}^d) \cong \Z_{p^{(s-t_1)}} \times \Z_{p^{(s-t_2)}} \times \cdots \times \Z_{p^{(s-t_j)}} \times (\Z_{p^s})^{\frac{1}{2}d(d-1)+(2k+1)-3j}.$$
$$M(G_{j,k}^d) \cong (\Z_{p^s})^{\binom{d}{2}+(2k+1)-3j} \times \prod_{i=1}^{j}\Z_{p^{(s-t_i)}} .
$$
\item [(ii)] If $j=k+1$, then 
%$$M(G_{j,k}^d) \cong \Z_{p^{(s-t_1)}} \times \Z_{p^{(s-t_2)}} \times \cdots \times \Z_{p^{(s-t_{k+1})}} \times (\Z_{p^s})^{\frac{1}{2}d(d-1)-(k+1)}.$$
$$M(G_{j,k}^d) \cong (\Z_{p^s})^{\binom{d}{2}-(k+1)} \times\prod_{i=1}^{k+1}\Z_{p^{(s-t_i)}} .$$
\end{enumerate}
\end{lemma}

\begin{theorem} \label{bound M(G_{j,k})}
 Let $d, r, s , m_i \in \mathbb{N}$ with $1 \leq r < d$ and
 %$r \leq j$, 
$0 < m_i < s$ for $i \in \{1, \ldots,r\}$.
 %, $0 \leq k \leq d-2$, ${1 \leq j \leq k+1}$
The following abelian groups occur as the Schur multiplier of some non-abelian group:
%of the group $G_{j,k}^d$ defined in Lemma \ref{SchurG_{j,k}}
 $$ (\Z_{p^s})^{\binom{d}{2}-(d-t)} \times  \Z_{p^{m_1}} \times \Z_{p^{m_2}} \times \cdots \times \Z_{p^{m_r}},$$
 where $t$ achieves all the following possible values
\begin{enumerate}
\item [(i)] $t=1$, if $r=d-1$.
    \item [(ii)]  $ t=1,2,3$, if $r=d-2$.
 \item [(iii)]  $1  \le t \le 3(d-r-1)$ except ${t=3(d-r-1)-1}$, if $r< d-2$.
\end{enumerate}
%for the pairs $(r,t) \in \{(d-1,1), (d-2,l), (m,n) \mid l=1,2,3, m<d-2, {1 \le n \le 3(d-r-1)}, n\neq 3(d-r-1)-1\}$
\end{theorem}

\begin{proof}
  Consider the group $G_{j,k}^d$ given in Lemma \ref{SchurG_{j,k}} with $r \le j$ and
 \[t_i=\biggl\{ 
\begin{array}{lcl}
s-m_i, && \text{ if } 1 \le i \le r\\
s, && \text{ if }r<i\le j.
	     \end{array}
\]
  %$t_i=s-m_i$ for $1 \le i \le r$, $t_i=s$ for $r<i\le j$ \sabnam{make cases}.
 Then we have,
$$M(G_{j,k}^d) \cong (\Z_{p^s})^{n_{j,k}} \times\prod_{i=1}^{r}\Z_{p^{m_i}} $$
for some $n_{j,k} \in \mathbb{N}$.  
  %Since $r \le j \le d-1$ and $(r-1) \le k \le (d-2)$, varying $j$ and $k$ yields a $(d-r) \times (d-r)$ square matrix, say $M$, whose elements are the number of copies of $\Z_{p^s}$ appearing in $M(G_{j,k}^d)$. 
 For $r \le j \le (d-1)$ and $(r-1) \le k \le (d-2)$, define a matrix,
 $$A=[m_{j,k}] \in M_{d-r}(\mathbb Z),$$
having entries
\[ m_{j,k}=\biggl\{ 
\begin{array}{lcl}
n_{j,k}, && \text{if } j \leq k+1\\
0, && \text{otherwise}.
	     \end{array}
\]
   %$m_{j,k} $ denotes the number of copies of $\Z_{p^s}$ appearing in  $M(G_{j,k}^d)$ when $j \le k+1$, and zero otherwise.
   %\textcolor{red}{To establish the result, it is enough to observe the entries of the principal diagonal and the diagonals that are parallel to it.} 
For $-1 \leq a \leq d-r-2$,   define a set ${S_a=\{m_{j,k}\mid k-j=a\}}$. %where \textcolor{red}{$-1 \leq a$}.
  %\textcolor{red}{The set $S_a$ contains the elements of the $(a+1)$-th diagonal entries of $M$.} 
  Note that when $a=-1$, it gives the principal diagonal entries of $A$. By Lemma \ref{SchurG_{j,k}}$(ii)$, ${\{m_{(k+1),k} \mid (r-1) \le k \le (d-2)\}}$ is the set of all integers between ${\min S_{-1}=\binom{d}{2}-(d-1)}$ and $\max S_{-1}={\binom{d}{2}-r}$. 
  
%  \begin{align}\label{values of m_jk}
%      \binom{d}{2}-(d-1) \le m_{(k+1),k} \le \binom{d}{2}-r.
%  \end{align}
If $r=d-1$, then $k=d-2$ and $m_{(d-1),(d-2)}=\binom{d}{2}-(d-1)$. Therefore,
$t=1$ and hence $\textit{(i)}$ follows.

For $r=d-2$, $k$ takes the values $d-3,d-2$. Using Lemma \ref{SchurG_{j,k}}, we have 
$$m_{(d-2),(d-3)}=\binom{d}{2}-(d-2),\quad m_{(d-1),(d-2)}=\binom{d}{2}-(d-1),$$
$$m_{(d-2),(d-2)}=\binom{d}{2}-(d-3).$$
Therefore, $t=1,2,3$ and hence $(ii)$ follows.

Now assume that $r<d-2$.  If $a=k-j$, then ${r \le j \le (d-2)-a}$ and it follows that ${a+r \le k \le d-2}$. 
Consider ${0 \le a \le (d-r-2)}$. By Lemma~\ref{SchurG_{j,k}}$(i)$, $m_{k-a,k}
%={\binom{d}{2}+(2k+1)-3(k-a)}
=\binom{d}{2}-k+(3a+1)$. Since $\min {S_{a}}=\binom{d}{2}+3a-(d-3) $ and $ {\max{S_{a}}= \binom{d}{2}+3a -(a+r-1)}$, the set
$S_a={\{m_{k-a,k}\mid a+r \le k \le d-2\}}$  consists of all the integers $n$ such that
$$ \binom{d}{2}+3a-(d-3) \leq n \leq  \binom{d}{2}+3a -(a+r-1).$$
%Note that $m_{jk}$ achieves the value $\binom{d}{2}+3a_0-(d-3)$ when $j=d-2-a_0$ and $k=d-2$. It achieves the value $\binom{d}{2}+3a_0-(a_0+r-1)$ when $j=r$ and $k=a_0+r$.
%For $a=a_0+1$, we have $r \le j \le (d-2)-(a_0+1)$ and $(a_0+1+r) \le k \le (d-2).$ Then again by Lemma \ref{SchurG_{j,k}}$(i)$, $m_{(k-a_0-1),k}%={\binom{d}{2}+(2k+1)-3(k-a_0-1)}
%=\binom{d}{2}-k+(3a_0+4).$
%Therefore,
%$$\binom{d}{2}+3a_0-(d-6) \le m_{jk} \le \binom{d}{2}+3a_0 -(a_0+r-3).$$
Also, for  $0\leq a \leq d-r-5$, we have $${\min{S_{a+1}}=\binom{d}{2}+3(a+1)-(d-3)\le \max{S_{a}}=\binom{d}{2}+3a -(a+r-1)}.$$ 
Thus, 
${\{m_{k-a,k} \mid 0\leq a \leq d-r-5,\ a+r \leq k \leq d-2\}}$  is the set of all integers between $\min{S_0}=\binom{d}{2}-(d-3)$ and ${\max{S_{d-r-5}}=\binom{d}{2}-(d-(3d-3r-9))}$. 
%for consecutive values of $a$ given by $a_0$ and $a_0+1$ we have $$\binom{d}{2}+3(a_0+1)-(d-3) \le \binom{d}{2}+3a_0 -(a_0+r-1) \implies a_0 \le d-r-5.$$
Hence, we are left with ${a=d-r-4}$, ${d-r-3}$ and $d-r-2$. The following table gives the values of $m_{j,k}$ for those $a$.
\[\setlength{\tabcolsep}{3pt} % Default value: 6pt
\renewcommand{\arraystretch}{2.2} % Default value: 
{\footnotesize
\begin{tabular}{|c|c |} 
	\hline
	$a$ &  $m_{j,k}$ \\ [0.005ex] 
	\hline\hline 
    $d-r-4$ & $\binom{d}{2}-(d-(3d-3r-7)),~\binom{d}{2}-(d-(3d-3r-8)),~\binom{d}{2}-(d-(3d-3r-9))$\\
    \hline
    $d-r-3$ & $\binom{d}{2}-(d-(3d-3r-5)),~\binom{d}{2}-(d-(3d-3r-6)$\\
    \hline
    $d-r-2$ & $\binom{d}{2}-(d-(3d-3r-3))$\\
    \hline
    \end{tabular}}\]\\
    % Now for $a=d-r-4$ we have $r \le j \le r+2$ and $d-4 \le k \le d-2.$ This gives
 %  $$m_{r,(d-4)}=\binom{d}{2}+2d-(3r+7),\quad m_{(r+1),(d-3)}=\binom{d}{2}+2d-(3r+8)$$
 %  $$m_{(r+2),(d-2)}=\binom{d}{2}+2d-(3r+9).$$
 %  Next for $a=d-r-3$, we have $j=r,r+1$ and $k=d-3,d-2.$ Then
 % $$m_{r,(d-3)}=\binom{d}{2}+2d-(3r+5), \quad m_{(r+1),(d-2)}=\binom{d}{2}+2d-(3r+6).$$
% And for $a=d-r-2$, $j=r$ and $k=d-2.$ So $m_{r,(d-2)}=\binom{d}{2}+2d-(3r+3).$
Therefore, ${\{m_{k-a,k} \mid 0\leq a \leq d-r-2,\ a+r \leq k \leq d-2\}}$ is the set of all integers lying between $\binom{d}{2}-(d-3)$ and 
 $\binom{d}{2}-(d-(3d-3r-3))$ except $\binom{d}{2}-(d-(3d-3r-4))$.
 Now for $a=-1$, we have
 $$m_{(d-2),(d-3)}=\binom{d}{2}-(d-2),\quad m_{(d-1),(d-2)}=\binom{d}{2}-(d-1).$$ 
 %Therefore, ${1 \le t \le 3(d-r-1)}$ except $t=3(d-r-1)-1.$
 %\begin{align*}
%     & \binom{d}{2}-(d-1) \le m_{jk} \le \binom{d}{2}+2d-(3r+3) \text{ except } \binom{d}{2}+2d-(3r+4)\\
%     \implies &  \binom{d}{2}-(d-1) \le m_{jk} \le \binom{d}{2}-\{d-(3d-3r-3)\}\\
% & \text{ except } \binom{d}{2}-\{d-(3d-3r-4)\}.
%  \end{align*}
This proves $(iii)$ and hence the result.
\end{proof} 

\end{document}
